\documentclass[11pt,a4paper]{article}
\usepackage[left=3cm,right=3cm,top=3cm,bottom=3cm]{geometry}
\usepackage{lmodern}
\usepackage{setspace}
\usepackage{parskip}
\usepackage{hyperref}
\usepackage{amsmath, amssymb}
\usepackage{graphicx}
% \usepackage{natbib}
% \usepackage[backend=biber, style=authoryear, natbib=true]{biblatex}
% \usepackage[backend=biber, style=apa, natbib=true]{biblatex}
\usepackage[backend=biber, style=authoryear, natbib=true, url=false, doi=false]{biblatex}
\addbibresource{references1.bib}

% \bibliographystyle{apalike}

\title{Modeling Strategic Decision-Making in Dynamic Markets: Research Interests and Contributions}
%  \\
\author{%
Zhongli Wang\thanks{PhD student affiliated with Bielefeld University (UNIBI) and Catholic University of Milan (UCSC). Email address: \url{zhongli.wang@uni-bielefeld.de}. Many thanks for the guidance of Prof. Herbert Dawid (UNIBI), Prof. Domenico Delli Gatti (UCSC) and Prof. Christiane Fuchs (UNIBI).} \\\\%
% \href{mailto:nkwangzhongli@gmail.com}{nkwangzhongli@gmail.com} \\\\%
% Bielefeld University%
}
% \date{\today}


\begin{document}

\maketitle

% \begin{abstract}
% This short essay outlines my core research interests in strategic management and organizational adaptation, and illustrates how these themes align with the strategic and methodological focus of the Chair of International Management. I reflect on past and ongoing research, methodological approaches, and future directions for collaboration.
% \end{abstract}

My research examines how organizations develop and deploy strategic capabilities in competitive environments, focusing on the dynamic interactions between firm strategy, organizational adaptation, and industry evolution. 
My academic background and doctoral research are grounded in these themes, with a particular focus on computational modeling approaches to strategic management questions.
In my dissertation, titled \textit{Strategic Behavior and Economic Dynamics in ICT Markets: An Agent-Based Framework}, I investigate how organizations develop and reconfigure their strategic capabilities over time in response to technological disruption and evolving stakeholder demand. 
This work integrates firm competition theories with computational modeling to analyze organizational decision-making under various scenarios.

My broader research interests include strategic adaptation and organizational learning in dynamic environments, platform strategies and ecosystem management, 
organizational heterogeneity and capability development, strategic implications of regulatory and policy environments, and digital transformation and strategic analytics in emerging business models.
Methodologically, I employ innovative computational approaches to examine organizational phenomena and strategic decision-making processes. 
My toolkit includes Agent-Based Modeling (ABM) for organizational interactions and strategic dynamics, network-based approaches to stakeholder interactions and ecosystem dynamics, and applied machine learning.
These methodologies facilitate the exploration of organizational heterogeneity, strategic learning processes, and non-linear competitive dynamics that characterize complex business environments.
Relating to potential institutional fit, the Chair's emphasis on strategic management, international research orientation, and methodological rigor suits my research profile. 
In particular, my experience in simulation methods and firm behavior under uncertainty aligns with ongoing work at the Chair.

I am especially interested in potential collaboration on strategic decision-making and capability development in global organizations, organizational adaptation and evolution in rapidly changing environments, 
competitive dynamics and strategic renewal through technological innovation, 
methodological advances in computational approaches to strategic management questions.
Building on this foundation, in my ongoing and future research, I am particularly interested in exploring several key questions aligned with contemporary strategic management challenges:

One key research direction examines how organizations in emerging technology sectors—such as robotics and large language models—orchestrate strategic renewal processes under uncertainty and performance pressure.
Building on \textcite{Keil2008ThePerformance}, I propose to develop an agent-based simulation model that examines how firms select governance modes for external innovation—corporate venture capital (CVC), alliances, joint ventures, or acquisitions—based on inter-organizational trust, technological relatedness, and environmental turbulence. 
At the same time, drawing on insights from \textcite{Posen2017RenewingAgenda}, I plan to integrate a behavioral logic of problemistic search, where firms only adapt or initiate governance changes when performance falls below aspiration levels. 
The simulation would capture how firms experiment, imitate, or renew their strategies in response to perceived underperformance, while facing evolving market dynamics and innovation pressure. 
Through the research, it is possible to check whether firms relying solely on static governance routines may miss windows of opportunity or overcommit to misaligned innovation modes. 
By combining the behavioral theory of the firm with strategy process research on governance and external venturing, this project would contribute to a deeper understanding of how firms reconfigure their innovation strategies through adaptive governance under uncertainty.

Another research area investigates the strategic decision-making processes behind organizational choices to limit international expansion.
This line of inquiry focuses on the strategic decision-making processes behind internationalization choices among technology-intensive organizations. 
While much of the strategic management literature emphasizes global growth trajectories \parencite{Cuervo-Cazurra2015InternationalizationEscape}, relatively less attention is paid to why some high-potential organizations deliberately choose to remain focused on domestic markets, especially in highly scalable sectors like digital platforms and AI. 
Drawing on the explanation of real option theory \parencite{Brouthers2008RealPerformance} and the attention-based view of the firm \parencite{March1991ExplorationLearning}, this project aims to explore how leadership teams interpret strategic signals and make internationalization decisions under complexity, ambiguity, and resource constraints. 
I plan to develop computational models that capture how strategic decision-makers evaluate foreign market opportunities based on imperfect information regarding stakeholder demands, institutional environments, and competitive dynamics. 
These models would allow examination of how managerial cognition, organizational identity, and strategic frames influence expansion choices. 
Through the research, it is possible to check whether organizations with substantial resources may strategically delay or limit international expansion due to attentional constraints, identity considerations, or strategic focus decisions. 
The research would contribute to the strategic management literature by reconceptualizing limited internationalization not as organizational failure, but as deliberate strategic positioning—with implications for leadership development and strategic decision processes in technology ventures.

Finally, I am interested in examining how global organizations in the AI sector develop non-market strategies and dynamic capabilities to navigate complex regulatory environments while creating stakeholder value across diverse user ecosystems.
As organizations like OpenAI, DeepSeek, and Mistral pursue global strategies, they must navigate complex stakeholder ecosystems, diverse user demands, and increasingly fragmented institutional environments. 
Emerging governance frameworks such as the EU AI Act, China's algorithm oversight mechanisms, and U.S. strategic technology policies are reshaping the competitive landscape, raising critical questions about organizational adaptation and strategic positioning in regulated innovation contexts. This project aims to explore how global organizations develop and deploy integrated market and non-market strategies—encompassing value creation, stakeholder engagement, and institutional navigation—in response to regulatory heterogeneity and evolving stakeholder expectations. Using computational modeling approaches combined with strategic management frameworks \parencite{Teece1997DynamicManagement, Barney1991FirmAdvantage}, I plan to investigate how organizations develop distinctive capabilities in environment-sensing, strategic flexibility, and cross-border knowledge integration to create sustainable competitive advantage in regulated innovation ecosystems \parencite{Parker2017PlatformEcosystems}. This research would contribute to strategic management literature on dynamic capabilities, corporate political activity, and organizational adaptation in complex institutional environments, while offering valuable insights into how technology organizations can align strategic decision processes with both stakeholder value creation imperatives and institutional legitimacy requirements.

Through these research directions, my goal is to further develop a rigorous research portfolio in strategic management, combining theory development, computational methodologies, and empirical relevance to organizational practice. I see strong opportunities for synergy with the Chair's focus areas, and I am eager to contribute to an ambitious and collaborative research environment.

% \bibliography{references}
\printbibliography
\end{document}

% \vspace{2em}
% \noindent\textbf{Keywords:} Industrial Organization, Agent-Based Modeling, Firm Strategy, Economic Dynamics, Simulation Methods

A third research direction explores the relationship between external stakeholder expectations and organizational strategy in the context of emerging technologies, examining specifically how resource provider expectations and capital ecosystem dynamics shape strategic adaptation and organizational resilience in nascent industries.
This area examines the strategic implications of resource availability cycles in the robotics sector. In recent years, numerous technology ventures in China and beyond have experienced rapid capability development supported by both institutional incentives and venture funding. 
However, some organizations subsequently faced abrupt resource constraints as stakeholder expectations shifted, necessitating strategic pivots, organizational restructuring, or market exit. 
This pattern raises important questions about how organizations develop resilient strategies under resource volatility \parencite{Gottfried2023InvestigationInvestors}, and how stakeholder-organization relationships shape competitive dynamics in emerging sectors. 
Using computational modeling approaches, I plan to investigate the co-evolution of organizational decision-making and resource provider behavior, exploring how expectation cycles, perform mance indicators, and institutional support mechanisms influence strategic renewal processes, capability development trajectories, and ecosystem formation. Preliminary findings suggest that transient resource availability can distort organizational priorities, leading to premature scaling and ecosystem fragility, whereas organizations with stable resource access or adaptive capability development approaches demonstrate greater resilience and long-term leadership. This research contributes to strategic management literature on resource orchestration \parencite{Sirmon2007ManagingBox}, organizational resilience, and innovation ecosystems, while offering insights into how governance structures and stakeholder relationships influence sustainable competitive advantage in emerging fields.
